\documentclass[11pt,a4paper]{article}

% =================================================
% Paquetes
% =================================================

\usepackage[utf8]{inputenc}
\usepackage[T1]{fontenc}
\usepackage[spanish]{babel}
\usepackage{geometry}
\usepackage{enumitem}
\usepackage{hyperref}
\usepackage{xcolor}
\usepackage{listings}
\usepackage{booktabs}
\usepackage{array}
\usepackage{tabularx}
\usepackage{parskip}

\geometry{
  top=2.5cm,
  bottom=2.5cm,
  left=2.7cm,
  right=2.7cm
}

\hypersetup{
  colorlinks=true,
  linkcolor=black,
  urlcolor=blue
}

\setlist[itemize]{
  noitemsep,
  topsep=4pt
}

\setlist[enumerate]{
  noitemsep,
  topsep=4pt
}

\lstset{
  basicstyle=\ttfamily\small,
  frame=single,
  breaklines=true,
  columns=fullflexible,
  showstringspaces=false
}

\lstset{
  literate=
  {á}{{\'a}}1 {é}{{\'e}}1 {í}{{\'i}}1 {ó}{{\'o}}1 {ú}{{\'u}}1
  {Á}{{\'A}}1 {É}{{\'E}}1 {Í}{{\'I}}1 {Ó}{{\'O}}1 {Ú}{{\'U}}1
  {ñ}{{\~n}}1 {Ñ}{{\~N}}1 {¿}{{?`}}1 {¡}{{!`}}1
}

% =================================================
% Documento
% =================================================

\title{
  \textbf{Guía de Tecnologías, Arquitectura y Despliegue}\\
  \large Taller de Integración II y Taller de Integración IV
}

\author{}
\date{Segundo semestre 2026}

\begin{document}

\maketitle

% =================================================
\section{Objetivo}
% =================================================

El presente documento define las principales tecnologías, la arquitectura
general, la estructura del monorepo y los procedimientos básicos de
desarrollo y despliegue utilizados durante el trabajo conjunto de Taller de
Integración II (TI2) y Taller de Integración IV (TI4).

Los objetivos de esta guía son:

\begin{itemize}
  \item establecer un stack tecnológico común;
  \item explicar la responsabilidad de cada tecnología;
  \item justificar las principales decisiones tecnológicas;
  \item definir la estructura inicial del monorepo;
  \item explicar cómo ejecutar los proyectos localmente;
  \item definir cómo se comunican las aplicaciones con el backend;
  \item establecer cómo se gestionan variables de entorno y secretos;
  \item explicar cómo desplegar las diferentes partes del sistema;
  \item establecer claramente las responsabilidades técnicas de TI2 y TI4;
  \item reducir decisiones técnicas innecesarias durante el desarrollo.
\end{itemize}

Este documento complementa el protocolo de trabajo con Git, GitHub y Linear.

% =================================================
\section{Principio general de arquitectura}
% =================================================

El proyecto corresponde a un único producto y se organizará en \textbf{un
monorepo GitHub}. El monorepo contendrá las aplicaciones Web y Mobile, el
backend compartido y los paquetes internos que se incorporen durante el
desarrollo.

\begin{lstlisting}
Proyecto
|
+-- apps/
|   +-- web/                 Aplicación Web de TI2
|   +-- mobile/              Aplicación Mobile de TI4
|
+-- convex/                  Backend, base de datos y autenticación
+-- packages/                Código compartido cuando corresponda
+-- .github/workflows/       Integración continua
+-- package.json             Configuración de workspaces
+-- bun.lock                 Único archivo de bloqueo
\end{lstlisting}

Ambas aplicaciones utilizarán:

\begin{itemize}
  \item el mismo backend;
  \item la misma base de datos;
  \item el mismo sistema de autenticación;
  \item las mismas reglas de negocio.
\end{itemize}

La arquitectura general será:

\begin{lstlisting}
                     CONVEX
              Backend + Database
                     |
                Better Auth
                /          \
               /            \
              v              v
        Web - TI2        Mobile - TI4

          React          React Native
           Vite              Expo
     TanStack Router      Expo Router
\end{lstlisting}

Por lo tanto, TI2 será propietario técnico del backend compartido y TI4
consumirá los servicios ofrecidos por TI2 desde la misma base de código.

% =================================================
\section{Stack tecnológico}
% =================================================

Las principales tecnologías seleccionadas son:

\begin{table}[h]
  \centering
  \begin{tabularx}{\textwidth}{>{\bfseries}l X}
    \toprule
    Componente & Tecnología \\
    \midrule
    Lenguaje principal & TypeScript \\
    Frontend Web & React \\
    Build Web & Vite \\
    Routing Web & TanStack Router \\
    Aplicación móvil & React Native + Expo \\
    Routing móvil & Expo Router \\
    Backend & Convex \\
    Base de Datos & Convex \\
    Autenticación & Better Auth \\
    Build móvil & EAS Build \\
    Hosting Web & Vercel \\
    Hosting Backend & Convex Cloud \\
    Control de versiones & Git + GitHub \\
    Gestión del trabajo & Linear \\
    Integración continua & GitHub Actions \\
    Organización del código & Monorepo con Bun workspaces \\
    \bottomrule
  \end{tabularx}
\end{table}

Los comandos de esta guía utilizarán \texttt{bun} como gestor de paquetes y se
ejecutarán desde la raíz del monorepo, salvo cuando se indique explícitamente
una aplicación.

Se deberá mantener un único gestor de paquetes y un único archivo de bloqueo en
la raíz del monorepo. Las aplicaciones y paquetes internos utilizarán los
workspaces definidos en el \texttt{package.json} raíz.

\subsection{Instalación de Bun}

Antes de ejecutar los comandos de esta guía, cada integrante deberá instalar
Bun en su equipo y verificar que esté disponible en la terminal.

En macOS y Linux se utilizará:

\begin{lstlisting}
curl -fsSL https://bun.sh/install | bash
\end{lstlisting}

En Windows se utilizará PowerShell:

\begin{lstlisting}
powershell -c "irm bun.sh/install.ps1|iex"
\end{lstlisting}

La mayoría de los integrantes de TI2 utilizará Windows, por lo que esa será la
instalación de referencia para dicho equipo. Después de instalar Bun, se deberá
abrir una nueva terminal y comprobar la instalación mediante:

\begin{lstlisting}
bun --version
\end{lstlisting}

Por ejemplo, no deberán coexistir en el monorepo:

\begin{lstlisting}
            bun.lock
            package-lock.json
            pnpm-lock.yaml
            yarn.lock
\end{lstlisting}

% =================================================
\section{TypeScript}
% =================================================

\subsection{Responsabilidad}

TypeScript será el lenguaje principal de programación de ambos equipos.

Será utilizado en:

\begin{itemize}
  \item frontend web;
  \item aplicación móvil;
  \item funciones del backend;
  \item definición del modelo de datos;
  \item autenticación;
  \item contratos entre clientes y backend.
\end{itemize}

\subsection{Justificación}

Utilizar TypeScript en todo el proyecto permite compartir un mismo modelo
mental entre ambos equipos y detectar una gran cantidad de errores durante el
desarrollo antes de ejecutar la aplicación.

Además, Convex permite obtener tipos de extremo a extremo entre sus funciones
de backend y los clientes TypeScript.

Como regla general, el código nuevo deberá escribirse en TypeScript.

% =================================================
\section{React}
% =================================================

React será utilizado para construir la interfaz de usuario de la aplicación
web de TI2.

React será responsable de:

\begin{itemize}
  \item componentes visuales;
  \item formularios;
  \item composición de pantallas;
  \item interacción con el usuario;
  \item presentación de los datos obtenidos desde Convex.
\end{itemize}

React no será responsable de acceder directamente a la base de datos.

El flujo esperado será:

\begin{lstlisting}
Componente React
      |
      v
Cliente Convex
      |
      v
Función Convex
      |
      v
Base de Datos
\end{lstlisting}

% =================================================
\section{Vite}
% =================================================

Vite será utilizado como herramienta de desarrollo y construcción de la
aplicación web.

Sus responsabilidades principales serán:

\begin{itemize}
  \item ejecutar el servidor de desarrollo;
  \item procesar TypeScript y JSX;
  \item administrar módulos durante desarrollo;
  \item generar el build de producción.
\end{itemize}

\subsection{Justificación}

La aplicación web funcionará principalmente como cliente de Convex.

Por esta razón no es necesario introducir un segundo framework full-stack con
una capa adicional de servidor.

La arquitectura será deliberadamente simple:

\begin{lstlisting}
React
 |
 +-- TanStack Router
 |
 +-- Convex Client
 |
 +-- Better Auth Client
\end{lstlisting}

La lógica de servidor permanecerá en Convex.

% =================================================
\section{TanStack Router}
% =================================================

TanStack Router será utilizado como sistema de routing de la aplicación web.

Será responsable de:

\begin{itemize}
  \item definición de páginas;
  \item navegación;
  \item parámetros de URL;
  \item layouts;
  \item rutas anidadas;
  \item navegación type-safe.
\end{itemize}

Se utilizará routing basado en archivos.

Por defecto las rutas se almacenarán en:

\begin{lstlisting}
src/routes/
\end{lstlisting}

TanStack Router generará automáticamente:

\begin{lstlisting}
src/routeTree.gen.ts
\end{lstlisting}

Este archivo es generado automáticamente y no deberá modificarse manualmente.

\subsection{Justificación}

Se utilizará TanStack Router en lugar de TanStack Start porque el servidor del
proyecto ya está definido mediante Convex.

TanStack Start incorporaría funcionalidades adicionales de servidor que no
son necesarias para la arquitectura definida.

% =================================================
\section{Convex}
% =================================================

Convex será la plataforma central del backend.

Será utilizado como:

\begin{itemize}
  \item backend;
  \item base de datos;
  \item plataforma de ejecución de funciones;
  \item mecanismo de sincronización de datos con los clientes.
\end{itemize}

El backend será desarrollado y mantenido por TI2.

\subsection{Tipos de funciones}

La lógica principal se implementará mediante:

\begin{description}
  \item[Queries:] operaciones de lectura.
  \item[Mutations:] operaciones que modifican datos.
  \item[Actions:] operaciones que requieren interacción con servicios externos
    o comportamientos que no corresponden a una query o mutation.
\end{description}

Conceptualmente:

\begin{lstlisting}
Cliente
  |
  +-- Query -------> Leer
  |
  +-- Mutation ----> Modificar
  |
  +-- Action ------> Servicio externo / proceso especial
\end{lstlisting}

La mayor parte de la lógica de negocio deberá permanecer en queries y
mutations cuando corresponda.

% =================================================
\section{Base de datos Convex}
% =================================================

El modelo de datos será administrado mediante Convex.

El esquema principal se encontrará en:

\begin{lstlisting}
convex/schema.ts
\end{lstlisting}

El esquema deberá formar parte del control de versiones.

Los cambios relevantes al modelo de datos deberán:

\begin{itemize}
  \item realizarse mediante una branch;
  \item pasar por Pull Request;
  \item considerar el impacto sobre la aplicación web;
  \item considerar el impacto sobre la aplicación móvil;
  \item quedar asociados al issue correspondiente en Linear.
\end{itemize}

No deberán existir cambios manuales indispensables que solamente conozca un
integrante del equipo.

% =================================================
\section{Better Auth}
% =================================================

Better Auth administrará la autenticación de los usuarios.

La instancia principal de Better Auth será ejecutada sobre el backend Convex.

Entre sus responsabilidades se encuentran:

\begin{itemize}
  \item registro de usuarios;
  \item inicio de sesión;
  \item cierre de sesión;
  \item manejo de sesiones;
  \item mecanismos adicionales de autenticación habilitados por el proyecto.
\end{itemize}

\subsection{Autenticación y autorización}

Es importante distinguir:

\begin{description}
  \item[Autenticación:] determina quién es el usuario.
  \item[Autorización:] determina qué puede hacer el usuario.
\end{description}

Por lo tanto:

\begin{lstlisting}
Better Auth
    |
    v
¿Quién es el usuario?
    |
    v
Convex
    |
    v
¿Puede ejecutar esta operación?
\end{lstlisting}

Las reglas de autorización propias del negocio deberán implementarse en el
backend y no confiar exclusivamente en que una pantalla esté oculta en el
frontend.

% =================================================
\section{Expo y React Native}
% =================================================

TI4 desarrollará la aplicación móvil utilizando React Native mediante Expo.

Expo será responsable de proporcionar:

\begin{itemize}
  \item entorno de desarrollo móvil;
  \item herramientas de ejecución;
  \item configuración de Android e iOS;
  \item integración con APIs nativas;
  \item proceso de build mediante EAS.
\end{itemize}

\subsection{Justificación}

Expo reduce la cantidad de configuración nativa que el equipo debe administrar
directamente y permite mantener la mayor parte del desarrollo en TypeScript y
React.

Esto permite que ambos equipos compartan tecnologías y conceptos.

\begin{center}
  \begin{tabular}{ll}
    \toprule
    TI2 Web & TI4 Mobile \\
    \midrule
    React & React Native \\
    TanStack Router & Expo Router \\
    TypeScript & TypeScript \\
    Convex & Convex \\
    Better Auth & Better Auth \\
    \bottomrule
  \end{tabular}
\end{center}

% =================================================
\section{Expo Router}
% =================================================

Expo Router será utilizado para administrar la navegación de la aplicación
móvil.

Las pantallas estarán representadas mediante archivos dentro de:

\begin{lstlisting}
app/
\end{lstlisting}

Por ejemplo:

\begin{lstlisting}
app/
|
+-- _layout.tsx
+-- index.tsx
+-- login.tsx
|
+-- usuarios/
    |
    +-- index.tsx
    +-- [id].tsx
\end{lstlisting}

Expo Router será responsable de:

\begin{itemize}
  \item navegación;
  \item layouts;
  \item parámetros;
  \item deep links;
  \item organización de pantallas.
\end{itemize}

La lógica de negocio no deberá implementarse dentro del router.

% =================================================
\section{Monorepo y un solo backend}
% =================================================

El backend Convex existirá solamente en la carpeta \texttt{convex/} de la raíz
del monorepo. TI4 \textbf{no deberá crear una segunda carpeta Convex ni copiar
funciones del backend}.

La arquitectura correcta será:

\begin{lstlisting}
Monorepo
|
 +-- convex/
 |   |
 |   +-- schema.ts
 |   +-- users.ts
 |   +-- auth...
|
+-- apps/
    |
    +-- web/
    +-- mobile/
\end{lstlisting}

De esta forma existe una única implementación del servidor y ambas
aplicaciones consumen la misma API pública.

% =================================================
\section{Tipos compartidos dentro del monorepo}
% =================================================

Convex genera una representación TypeScript de su API pública. En el monorepo
esa representación estará disponible para las aplicaciones Web y Mobile desde
la misma base de código.

Conceptualmente:

\begin{lstlisting}
convex/
     |
     | define funciones y genera tipos
     v
convex/_generated/api
     |                 |
     v                 v
apps/web          apps/mobile
\end{lstlisting}

Ambas aplicaciones utilizarán los tipos generados por Convex mediante un alias
o un paquete interno definido por el monorepo. La configuración concreta del
alias podrá variar, pero no deberá duplicar la implementación del backend ni
mantener un contrato generado de forma independiente para cada aplicación.

La dependencia de Convex se declarará en los workspaces que corresponda y se
instalará desde la raíz:

\begin{lstlisting}
bun add convex --cwd apps/web
bun add convex --cwd apps/mobile
\end{lstlisting}

El comando de desarrollo del backend se ejecutará desde la raíz:

\begin{lstlisting}
bunx convex dev
\end{lstlisting}

La carpeta \texttt{convex/\_generated/} deberá mantenerse de acuerdo con la
política del proyecto para archivos generados. Si la configuración de Convex
permite regenerarla automáticamente en cada entorno, no se deberá crear una
segunda copia manual dentro de \texttt{apps/mobile/}.

\subsection{Regla}

Web y Mobile consumen la API pública definida en \texttt{convex/}.

TI4 puede consumir las funciones del backend y utilizar sus tipos,
pero no debe modificar
directamente las funciones Convex que no estén dentro de su responsabilidad.

% =================================================
\section{Estructura del monorepo}
% =================================================

La estructura inicial será aproximadamente:

\begin{lstlisting}
proyecto/
|
 +-- apps/
 |   |
 |   +-- web/
 |   |   +-- src/
 |   |   +-- public/
 |   |   +-- package.json
 |   |   +-- vite.config.ts
 |   |
 |   +-- mobile/
 |       +-- app/
 |       +-- components/
 |       +-- features/
 |       +-- package.json
 |       +-- app.json
|
 +-- packages/
 |   +-- ui/                 Componentes compartidos si corresponde
 |   +-- shared/             Tipos o utilidades compartidas si corresponde
|
 +-- convex/
 |   +-- _generated/
 |   +-- betterAuth/
 |   +-- schema.ts
 |   +-- auth.config.ts
 |   +-- convex.config.ts
 |   +-- http.ts
 |   +-- ...
|
 +-- .github/
 |   +-- workflows/
 |
 +-- .env.example
 +-- .gitignore
 +-- package.json
 +-- bun.lock
 +-- tsconfig.json
 +-- vercel.json
 +-- README.md
\end{lstlisting}

Los directorios \texttt{packages/ui/} y \texttt{packages/shared/} son
opcionales. Solo deberán crearse cuando exista una responsabilidad compartida
real; no se incorporarán como capas vacías por anticipación.

El directorio del backend:

\begin{lstlisting}
convex/_generated/
\end{lstlisting}

se generará desde la configuración de Convex y no se deberá modificar
manualmente.

% =================================================
La aplicación móvil utilizará el backend y los paquetes compartidos desde este
monorepo; no tendrá una copia local del backend ni un archivo de contrato
externo.

% =================================================
\section{Setup inicial del monorepo}
% =================================================

Todos los comandos siguientes se ejecutarán desde la raíz del monorepo.

\subsection{1. Crear la raíz y los workspaces}

La raíz deberá declarar los workspaces de las aplicaciones y de los paquetes
internos:

\begin{lstlisting}
{
  "name": "proyecto",
  "private": true,
  "workspaces": [
    "apps/*",
    "packages/*"
  ]
}
\end{lstlisting}

\subsection{2. Crear aplicación Web}

\begin{lstlisting}
bun create vite@latest apps/web -- --template react-ts
bun install
\end{lstlisting}

\subsection{3. Crear aplicación Mobile}

\begin{lstlisting}
bunx create-expo-app@latest apps/mobile
bun install
\end{lstlisting}

La plantilla estándar de Expo incorpora Expo Router.

\subsection{4. Instalar Convex en la raíz}

El backend Convex se configurará en la raíz del monorepo:

\begin{lstlisting}
bun add convex
bunx convex dev
\end{lstlisting}

La primera ejecución permitirá vincular o crear el proyecto Convex y generará
la carpeta:

\begin{lstlisting}
convex/
\end{lstlisting}

\subsection{5. Instalar TanStack Router en Web}

\begin{lstlisting}
bun add @tanstack/react-router --cwd apps/web
bun add -d @tanstack/router-plugin --cwd apps/web
bun add -d @tanstack/react-router-devtools --cwd apps/web
\end{lstlisting}

En Vite se deberá agregar el plugin de TanStack Router antes del plugin de
React.

Ejemplo:

\begin{lstlisting}
import { defineConfig } from 'vite'
import react from '@vitejs/plugin-react'
import { tanstackRouter } from '@tanstack/router-plugin/vite'

export default defineConfig({
  plugins: [
    tanstackRouter({
      target: 'react',
      autoCodeSplitting: true,
    }),
    react(),
  ],
})
\end{lstlisting}

\subsection{6. Instalar Better Auth en el backend}

\begin{lstlisting}
bun add better-auth @convex-dev/better-auth
\end{lstlisting}

La integración deberá realizarse dentro del backend Convex siguiendo la
estructura oficial de Better Auth para Convex.

Entre otros archivos, se utilizarán:

\begin{lstlisting}
convex/
|
+-- auth.config.ts
+-- convex.config.ts
+-- http.ts
|
+-- betterAuth/
    |
    +-- convex.config.ts
    +-- auth.ts
    +-- schema.ts
    +-- adapter.ts
\end{lstlisting}

El esquema de Better Auth deberá generarse mediante sus herramientas y no
mantenerse manualmente si puede evitarse.

% =================================================
\section{Dependencias de la aplicación Mobile}
% =================================================

TI4 solamente utilizará el cliente Convex y los tipos generados desde la carpeta
compartida del backend. No deberá crear ni desplegar una segunda implementación.

\subsection{1. Instalar Better Auth para Expo}

\begin{lstlisting}
bun add better-auth @better-auth/expo --cwd apps/mobile
bun add expo-network expo-secure-store --cwd apps/mobile
\end{lstlisting}

Better Auth utilizará Secure Store para almacenar información de sesión de
forma apropiada para la aplicación móvil.

Si el proyecto utiliza proveedores sociales, podrán requerirse además
dependencias propias de Expo como linking o web browser.

% =================================================
\section{Conexión de la aplicación Web a Convex}
% =================================================

La aplicación web utilizará una instancia de:

\begin{lstlisting}
ConvexReactClient
\end{lstlisting}

configurada mediante la URL del deployment.

La variable pública utilizada por Vite será:

\begin{lstlisting}
VITE_CONVEX_URL=
\end{lstlisting}

Dentro del frontend:

\begin{lstlisting}
const convex = new ConvexReactClient(
  import.meta.env.VITE_CONVEX_URL
)
\end{lstlisting}

El cliente deberá proveerse en la raíz de la aplicación.

% =================================================
\section{Conexión de la aplicación Mobile a Convex}
% =================================================

Expo utilizará el mismo cliente React de Convex.

La variable será:

\begin{lstlisting}
EXPO_PUBLIC_CONVEX_URL=
\end{lstlisting}

En la raíz de Expo Router podrá configurarse:

\begin{lstlisting}
const convex = new ConvexReactClient(
  process.env.EXPO_PUBLIC_CONVEX_URL!
)
\end{lstlisting}

y proveerse a la aplicación mediante el provider correspondiente.

Esto significa que:

\begin{lstlisting}
Web
  |
  +----\
        \
         > mismo deployment Convex
        /
Mobile /
\end{lstlisting}

% =================================================
\section{URLs de Convex}
% =================================================

Convex utiliza URLs con responsabilidades distintas.

Conceptualmente tendremos:

\begin{description}
  \item[Convex URL:] conexión del cliente React con las funciones Convex.
  \item[Convex Site URL:] endpoints HTTP, incluyendo rutas utilizadas por
    Better Auth cuando corresponda.
\end{description}

Por ejemplo, los clientes podrán necesitar variables equivalentes a:

\begin{lstlisting}
VITE_CONVEX_URL=
VITE_CONVEX_SITE_URL=
\end{lstlisting}

para web, y:

\begin{lstlisting}
EXPO_PUBLIC_CONVEX_URL=
EXPO_PUBLIC_CONVEX_SITE_URL=
\end{lstlisting}

para mobile.

Estas URLs permiten localizar servicios públicos.

\textbf{No son el lugar donde deben almacenarse secretos.}

% =================================================
\section{Variables de entorno}
% =================================================

Se distinguirán claramente variables públicas y secretos.

\subsection{Frontend Web}

Vite expone al código cliente variables con prefijo:

\begin{lstlisting}
VITE_
\end{lstlisting}

Por ejemplo:

\begin{lstlisting}
VITE_CONVEX_URL=
\end{lstlisting}

Toda variable \texttt{VITE\_} debe considerarse visible para el usuario final.

\subsection{Aplicación móvil}

Expo expone al código cliente las variables con prefijo:

\begin{lstlisting}
EXPO_PUBLIC_
\end{lstlisting}

Por ejemplo:

\begin{lstlisting}
EXPO_PUBLIC_CONVEX_URL=
\end{lstlisting}

Toda variable \texttt{EXPO\_PUBLIC\_} debe considerarse pública.

\subsection{Backend}

Los secretos utilizados por las funciones Convex deberán configurarse en el
entorno de Convex.

Por ejemplo:

\begin{lstlisting}
BETTER_AUTH_SECRET
\end{lstlisting}

Estos valores no deberán estar:

\begin{itemize}
  \item dentro del código;
  \item dentro del frontend;
  \item dentro de variables públicas;
  \item versionados en Git.
\end{itemize}

% =================================================
\section{Archivos de entorno}
% =================================================

Los archivos locales que contengan valores específicos de cada desarrollador no
deberán subirse al monorepo.

Cada aplicación podrá mantener su propio archivo local dentro de su directorio:

\begin{lstlisting}
apps/web/.env.local
apps/mobile/.env.local
\end{lstlisting}

El monorepo deberá mantener ejemplos de configuración sin secretos:

\begin{lstlisting}
apps/web/.env.example
apps/mobile/.env.example
\end{lstlisting}

con las variables necesarias pero sin secretos.

Ejemplo TI2:

\begin{lstlisting}
VITE_CONVEX_URL=
VITE_CONVEX_SITE_URL=
\end{lstlisting}

Ejemplo TI4:

\begin{lstlisting}
EXPO_PUBLIC_CONVEX_URL=
EXPO_PUBLIC_CONVEX_SITE_URL=
\end{lstlisting}

% =================================================
\section{Ejecución local del monorepo}
% =================================================

Durante el desarrollo normalmente se utilizarán al menos dos procesos desde la
raíz del monorepo.

\subsection{Terminal 1: Convex}

\begin{lstlisting}
bunx convex dev
\end{lstlisting}

Este proceso sincroniza los cambios del backend con el deployment de desarrollo.

\subsection{Terminal 2: Web}

\begin{lstlisting}
bun --cwd apps/web run dev
\end{lstlisting}

Conceptualmente:

\begin{lstlisting}
Terminal 1                 Terminal 2

bunx convex dev            bun --cwd apps/web run dev
      |                          |
      v                          v
Backend Convex             React + Vite
\end{lstlisting}

% =================================================
\section{Ejecución local de Mobile}
% =================================================

La aplicación móvil se ejecutará mediante Expo.

\begin{lstlisting}
bun --cwd apps/mobile run start
\end{lstlisting}

o equivalentemente:

\begin{lstlisting}
cd apps/mobile && bunx expo start
\end{lstlisting}

Desde esta sesión se podrá ejecutar la aplicación en:

\begin{itemize}
  \item Android;
  \item iOS cuando el entorno lo permita;
  \item dispositivo físico;
  \item emulador o simulador.
\end{itemize}

La aplicación Mobile deberá conectarse al deployment de desarrollo de Convex
definido para el proyecto.

No es necesario ejecutar el backend para modificar exclusivamente la aplicación
Mobile, aunque sí deberá utilizarse un deployment válido para probar la
integración.

% =================================================
\section{Ambientes}
% =================================================

Se distinguirán al menos:

\begin{description}
  \item[Development:] programación y pruebas durante el desarrollo.
  \item[Production:] versión estable utilizada para las demostraciones y
    entregas.
\end{description}

En móvil también podrá existir un ambiente:

\begin{description}
  \item[Preview:] build destinado a pruebas del equipo.
\end{description}

Los entornos deberán utilizar configuraciones separadas cuando sea necesario.

% =================================================
\section{Build de la aplicación Web}
% =================================================

Antes de considerar una versión lista para integración deberá poder generarse
el build de producción.

\begin{lstlisting}
bun --cwd apps/web run build
\end{lstlisting}

Una aplicación que funciona solamente con:

\begin{lstlisting}
bun --cwd apps/web run dev
\end{lstlisting}

pero cuyo build falla no se considera lista para integrar.

% =================================================
\section{Despliegue del backend Convex}
% =================================================

Durante desarrollo se utilizará:

\begin{lstlisting}
bunx convex dev
\end{lstlisting}

Para desplegar funciones, esquema e índices hacia producción se utilizará:

\begin{lstlisting}
bunx convex deploy
\end{lstlisting}

El deployment de producción será responsabilidad de TI2 y se ejecutará desde la
raíz del monorepo.

No todos los integrantes deberán manipular manualmente producción.

% =================================================
\section{Despliegue de la aplicación Web}
% =================================================

La aplicación web será desplegada mediante Vercel.

Vercel se conectará al único repositorio GitHub y se configurará con
\texttt{apps/web/} como directorio de la aplicación. Las órdenes de instalación
y build deberán respetar el \texttt{bun.lock} ubicado en la raíz.

El flujo general será:

\begin{lstlisting}
GitHub
    |
    v
   main
    |
    v
 Vercel
    |
    +---- Build apps/web
    |
    +---- Deploy Convex
    |
    v
Producción
\end{lstlisting}

Para integrar el despliegue de Convex con Vercel deberá configurarse un deploy
key de producción como variable segura en Vercel:

\begin{lstlisting}
CONVEX_DEPLOY_KEY
\end{lstlisting}

El build podrá ejecutar:

\begin{lstlisting}
bunx convex deploy --cmd 'bun --cwd apps/web run build'
\end{lstlisting}

De esta manera se reduce el riesgo de publicar un frontend incompatible con las
funciones Convex correspondientes.

% =================================================
\section{Configuración SPA en Vercel}
% =================================================

TanStack Router administrará las rutas en el cliente.

Por ello, al solicitar directamente una ruta como:

\begin{lstlisting}
/usuarios/123
\end{lstlisting}

el hosting deberá entregar la aplicación y permitir que TanStack Router
resuelva la ruta.

El repositorio podrá incluir:

\begin{lstlisting}
vercel.json
\end{lstlisting}

con una regla equivalente a:

\begin{lstlisting}
{
  "rewrites": [
    {
      "source": "/(.*)",
      "destination": "/index.html"
    }
  ]
}
\end{lstlisting}

% =================================================
\section{Build de la aplicación móvil}
% =================================================

Los builds para instalar serán generados mediante EAS Build. La configuración de
EAS deberá utilizar \texttt{apps/mobile/} como directorio de la aplicación y
compartir el lockfile de la raíz cuando el flujo de CI se ejecute desde el
monorepo.

Primero deberá configurarse EAS para el proyecto.

Una configuración simplificada de:

\begin{lstlisting}
eas.json
\end{lstlisting}

podrá incluir:

\begin{lstlisting}
{
  "build": {
    "development": {
      "developmentClient": true
    },
    "preview": {
      "distribution": "internal"
    },
    "production": {}
  }
}
\end{lstlisting}

% =================================================
\section{Build Preview de Mobile}
% =================================================

Durante el curso no será necesario publicar cada versión en Google Play o App
Store.

Para pruebas se utilizará preferentemente un build interno.

Android:

\begin{lstlisting}
eas build --platform android --profile preview
\end{lstlisting}

iOS:

\begin{lstlisting}
eas build --platform ios --profile preview
\end{lstlisting}

EAS proporcionará un build que podrá compartirse con integrantes y personas
autorizadas para pruebas.

En Android, este mecanismo permite disponer de un APK para instalar sin depender
de una publicación en Play Store.

% =================================================
\section{Build de producción Mobile}
% =================================================

Cuando corresponda generar una versión de producción:

\begin{lstlisting}
eas build --platform android
\end{lstlisting}

o:

\begin{lstlisting}
eas build --platform ios
\end{lstlisting}

También puede utilizarse:

\begin{lstlisting}
eas build --platform all
\end{lstlisting}

La publicación efectiva en tiendas será una decisión independiente del proceso
de desarrollo.

% =================================================
\section{Flujos de despliegue}
% =================================================

\subsection{TI2}

\begin{lstlisting}
Monorepo GitHub
  |
  v
main
  |
  v
Vercel
  |
  +--> React/Vite
  |
  +--> Convex Production
  |
  v
Web desplegada
\end{lstlisting}

\subsection{TI4}

\begin{lstlisting}
Monorepo GitHub
  |
  v
main
  |
  v
EAS Build
  |
  +--> Preview
  |
  +--> Production
  |
  v
Aplicación para instalar
\end{lstlisting}

% =================================================
\section{Integración Continua}
% =================================================

GitHub Actions verificará los cambios del monorepo antes de integrarlos.

La instalación deberá realizarse desde la raíz para utilizar el lockfile único:

\begin{lstlisting}
bun install --frozen-lockfile
\end{lstlisting}

Cuando los scripts correspondientes existan, deberán ejecutarse como mínimo las
verificaciones de los componentes afectados:

\begin{enumerate}
  \item instalación de dependencias desde la raíz;
  \item chequeo de TypeScript;
  \item lint;
  \item pruebas automatizadas disponibles;
  \item build de Web, Mobile o backend según las rutas modificadas.
\end{enumerate}

Conceptualmente:

\begin{lstlisting}
Pull Request
     |
     v
bun install --frozen-lockfile
     |
     v
Type Check
     |
     v
Lint
     |
     v
Tests
     |
     v
Build
     |
     v
Code Review
     |
     v
Merge
\end{lstlisting}

El despliegue no reemplaza al proceso de CI.

% =================================================
\section{Responsabilidades de TI2}
% =================================================

TI2 será responsable técnicamente de:

\begin{itemize}
  \item desarrollar la aplicación web;
  \item mantener React y TanStack Router;
  \item implementar las funciones Convex;
  \item definir y mantener el modelo de datos;
  \item implementar Better Auth;
  \item implementar reglas de autorización;
  \item desplegar el backend;
  \item desplegar la aplicación web;
  \item mantener disponibles los servicios requeridos por TI4;
  \item mantener estable la API pública consumida por Mobile;
  \item informar cambios incompatibles;
  \item mantener actualizados los tipos generados y los paquetes compartidos
    cuando corresponda.
\end{itemize}

TI2 es propietario técnico del backend compartido.

% =================================================
\section{Responsabilidades de TI4}
% =================================================

TI4 será responsable técnicamente de:

\begin{itemize}
  \item desarrollar la aplicación móvil mediante Expo;
  \item organizar la navegación mediante Expo Router;
  \item integrar la aplicación con Convex;
  \item integrar los flujos de Better Auth;
  \item consumir los tipos generados desde la configuración compartida;
  \item detectar y reportar problemas de integración;
  \item generar builds mediante EAS;
  \item evitar duplicar lógica perteneciente al backend.
\end{itemize}

Además, TI4 mantiene las responsabilidades de gestión, coordinación, seguimiento
y control de calidad definidas en el documento metodológico general.

% =================================================
\section{Frontera entre TI2 y TI4}
% =================================================

La frontera técnica entre ambos equipos estará formada principalmente por las
funciones públicas de Convex, el sistema de autenticación compartido y los
paquetes internos que sean consumidos por más de una aplicación.

TI4 no deberá:

\begin{itemize}
  \item crear una segunda base de datos;
  \item crear una segunda implementación del backend;
  \item copiar las funciones Convex dentro de \texttt{apps/mobile/};
  \item modificar directamente producción de TI2;
  \item duplicar reglas de negocio únicamente para hacer funcionar Mobile.
\end{itemize}

TI2 deberá mantener suficientemente documentadas las funciones públicas
consumidas por Mobile.

Como mínimo deberá conocerse:

\begin{itemize}
  \item propósito;
  \item argumentos;
  \item resultado;
  \item requerimientos de autenticación;
  \item errores relevantes.
\end{itemize}

% =================================================
\section{Cambios que afectan a ambos equipos}
% =================================================

Un cambio deberá coordinarse mediante Linear cuando modifique:

\begin{itemize}
  \item nombre de una función pública;
  \item argumentos de una función;
  \item tipo de retorno;
  \item reglas de autenticación;
  \item reglas de autorización;
  \item estructura de datos utilizada por ambos clientes;
  \item comportamiento funcional esperado por la aplicación Mobile;
  \item un paquete compartido utilizado por Web y Mobile.
\end{itemize}

Estas dependencias deberán registrarse en Linear.

Por ejemplo:

\begin{lstlisting}
TI2-42 Cambiar flujo de autenticación
              |
              +-- bloquea -->
                              TI4-18
                              Adaptar login mobile
\end{lstlisting}

% =================================================
\section{Qué corresponde a cada capa}
% =================================================

\subsection*{React}

Interfaz de usuario Web.

No contiene acceso directo a la base de datos.

\subsection*{TanStack Router}

Navegación Web.

No contiene lógica de negocio.

\subsection*{Expo / React Native}

Interfaz de usuario Mobile.

No reemplaza el backend.

\subsection*{Expo Router}

Navegación Mobile.

No contiene reglas de negocio compartidas.

\subsection*{Convex}

Backend, persistencia y reglas de negocio compartidas.

\subsection*{Better Auth}

Identidad y sesiones.

No reemplaza las reglas de autorización del sistema.

% =================================================
\section{Ejemplo de distribución de responsabilidades}
% =================================================

Supongamos que debe implementarse:

\begin{center}
  \textit{``Mostrar el perfil de un usuario autenticado.''}
\end{center}

La responsabilidad se distribuirá así:

\begin{lstlisting}
Better Auth
    |
    +--> identifica al usuario actual
    |
    v
Convex
    |
    +--> valida permisos
    +--> obtiene información
    |
    +-------------+
    |             |
    v             v
   Web          Mobile
  React       React Native
\end{lstlisting}

El frontend decide \textbf{cómo mostrar} la información.

El backend decide \textbf{qué información puede obtener} el usuario.

% =================================================
\section{Proceso para una funcionalidad compartida}
% =================================================

Una funcionalidad que afecte a ambos equipos podrá seguir el siguiente flujo:

\begin{enumerate}
  \item TI4 realiza el triage y se identifican las tareas necesarias.
  \item TI2 implementa o modifica las funciones necesarias en Convex.
  \item TI2 y TI4 integran el cambio mediante uno o más Pull Requests
    relacionados.
  \item se regeneran o actualizan los tipos públicos de Convex;
  \item TI4 implementa el consumo desde la aplicación Mobile;
  \item ambos equipos prueban la integración;
  \item los cambios llegan a los procesos de despliegue de Web y Mobile.
\end{enumerate}

Conceptualmente:

\begin{lstlisting}
Linear
   |
   +-------------------+
   |                   |
   v                   v
Issue TI2          Issue TI4
   |                   |
   v                   |
Backend Convex         |
   |                   |
   v                   |
Tipos generados -------+
                       |
                       v
                  apps/mobile
\end{lstlisting}

% =================================================
\section{Principios técnicos}
% =================================================

Durante el desarrollo se aplicarán los siguientes principios:

\begin{enumerate}
  \item existe un solo backend;
  \item existe una sola fuente de verdad para los datos;
  \item Web y Mobile consumen los mismos servicios;
  \item los secretos nunca se almacenan en código cliente;
  \item las reglas de autorización se validan en el backend;
  \item la API pública y los paquetes compartidos deben ser explícitos;
  \item los cambios incompatibles deben coordinarse antes de integrarse;
  \item el código generado necesario para compilar debe mantenerse
    correctamente versionado;
  \item una aplicación que no puede generar su build no está terminada;
  \item los entornos de desarrollo y producción deben mantenerse separados.
\end{enumerate}

% =================================================
\section{Resumen de arquitectura}
% =================================================

El stack tecnológico puede resumirse de la siguiente manera:

\begin{lstlisting}
                         Convex
                Backend + Base de Datos
                         |
                    Better Auth
                         |
             +-----------+-----------+
             |                       |
             v                       v

          apps/web                 apps/mobile

           React                  React Native
           Vite                       Expo
     TanStack Router              Expo Router
             |                       |
             +-----------+-----------+
                         |
                     TypeScript
\end{lstlisting}

El flujo de desarrollo y despliegue será:

\begin{lstlisting}
                       Linear
                         |
                         v
                      Issue
                         |
                         v
                  Feature Branch
                         |
                         v
                   Pull Request
                         |
                 +-------+-------+
                 |               |
                 v               v
                CI          Code Review
                 |               |
                 +-------+-------+
                         |
                         v
                        main
                         |
             +-----------+-----------+
             |                       |
             v                       v
         TI2 Deploy              TI4 Build
             |                       |
       Vercel + Convex            EAS Build
             |                       |
             v                       v
            Web                  Mobile App
\end{lstlisting}

% =================================================
\section{Referencias oficiales}
% =================================================

Las configuraciones específicas pueden evolucionar con nuevas versiones de las
herramientas. Antes de realizar cambios importantes deberán consultarse sus
documentaciones oficiales.

\begin{itemize}

  \item TanStack Router:
    \url{https://tanstack.com/router/latest}

  \item TanStack Router con Vite:
    \href{https://tanstack.com/router/latest/docs/installation/with-vite}{Documentación
    de TanStack Router con Vite}

  \item Convex:
    \url{https://docs.convex.dev/}

  \item Convex con React:
    \url{https://docs.convex.dev/client/react/overview}

  \item Convex con React Native:
    \url{https://docs.convex.dev/quickstart/react-native}

  \item Convex en Vercel:
    \url{https://docs.convex.dev/production/hosting/vercel}

  \item Better Auth:
    \url{https://better-auth.com/}

  \item Better Auth con Convex:
    \url{https://better-auth.com/docs/integrations/convex}

  \item Better Auth con Expo:
    \url{https://better-auth.com/docs/integrations/expo}

  \item Expo:
    \url{https://docs.expo.dev/}

  \item EAS Build:
    \url{https://docs.expo.dev/build/introduction/}

  \item EAS Internal Distribution:
    \url{https://docs.expo.dev/build/internal-distribution/}

\end{itemize}

\end{document}
